\lab{12}{Прыжки по индексу}
\task
Добавить ускоренное продвижение по постинг-листам. Подобрать шаг
на обучающих запросах, проверить на других запросах, исследовать
скорость и накладной размер, убедиться в неизменности выдачи.

\method
Постинг-лист делится на блоки по \(S\) документов. Для каждого блока
хранится запись размером 24 байта: последний \(docID\) (4 байта),
предыдущий \(docID\) перед началом блока (4), смещение блока (8),
порядковый номер первого постинга (4), число постингов блока (4).
Таблица позволяет искать первый блок, последний документ которого
не меньше целевого \(docID\), бинарным поиском.

После перехода курсор читает нужный блок последовательно. Предыдущий
\(docID\) необходим для восстановления абсолютных значений из разностей:
без него произвольный вход в середину сжатого потока был бы невозможен.
Проверка порядкового номера и границ защищает от неправильного смещения.
Шаг 0 отключает таблицу и служит базовой линией.

Малый шаг сокращает последовательное чтение, но увеличивает таблицу
и число поисков по ней. Для \(L\) постингов доступ к блоку требует
порядка \(\log(L/S)\) обращений и до \(S\) последовательных чтений.
Эта оценка не включает стоимость позиционных списков и дисковых
переходов. Для коротких списков размер служебной записи может
оказаться сопоставимым с самими данными.

\journal
В индекс добавлены таблицы прыжков и метод \code{Cursor::advance}.
Для эксперимента проверяются шаги 0, 4, 16, 64, 128 и 256.
Подготовлены 15 обучающих и 15 тестовых запросов. Каждый набор
проходит один прогрев и 15 измеряемых повторов. Отдельный инструмент
меняет таблицу шага для существующих постингов, что позволяет
сравнивать один корпус и одну нормализацию. Он читает готовый VB-индекс
без прыжков и создаёт вариант с заданным положительным шагом,
сохраняя поток постингов. Результат преобразования побайтово
сравнивается с обычным построением на тестовом корпусе.

Шаг выбирается минимумом среднего времени на train. Тестовые данные
используются после фиксации выбора. Поэтому вариант, у которого
случайно лучше медиана на test, не назначается победителем задним числом.
Подготовка шести вариантов выполняется отдельной командой
\code{just prepare-tune}. Команда \code{just tune} только измеряет
готовые индексы. До записи отчётов она проверяет наличие всех
вариантов, кодек, нормализацию и шаг. Исходный корпус для замеров
не нужен, существующие индексы не перестраиваются.

\testing
Тесты проверяют переход к существующему и отсутствующему документу,
за конец списка, к границе блока и в середину VB-потока. Одни и те же
булевы запросы и цитаты выполняются с прыжками и без них на raw,
VB и многоблочном индексе. Четыре модульных теста повторно пройдены.
В полном опыте контрольные суммы должны совпадать отдельно внутри
train и test для всех шагов.

\results
\begin{table}[H]\centering
\begin{tabular}{rrrrr}\toprule
Шаг & train, мс & test, мс & Постингов test & Прыжки, MiB\\\midrule
0 & 67,316 & 31,027 & 28\,832\,670 & 0,00\\
4 & 181,180 & 75,950 & 12\,635\,580 & 3975,52\\
16 & 166,591 & 71,556 & 18\,184\,740 & 1332,07\\
64 & 153,841 & 66,239 & 23\,225\,865 & 679,69\\
128 & 147,579 & 63,039 & 24\,700\,275 & 572,43\\
256 & 139,670 & 59,450 & 25\,719\,630 & 519,30\\
\bottomrule\end{tabular}
\caption{Время, суммарное чтение за 15 повторов и размер таблиц прыжков.}
\end{table}

\begin{figure}[H]\centering
\begin{tikzpicture}
\begin{axis}[width=0.95\textwidth,height=6.8cm,
xlabel={Шаг прыжков; 0 - без прыжков},ylabel={Среднее время, мс},
xtick={0,1,2,3,4,5},xticklabels={0,4,16,64,128,256},grid=major,
legend style={at={(0.98,0.98)},anchor=north east}]
\addplot[blue,mark=*] table[x=i,y=train]{figures/skips.dat};
\addlegendentry{train}
\addplot[red,mark=square*] table[x=i,y=test]{figures/skips.dat};
\addlegendentry{test}
\end{axis}
\end{tikzpicture}
\caption{Среднее время поиска при 15 повторах после прогрева.}
\end{figure}
На train лучший шаг - 0. Следовательно, выбранный вариант на test
совпадает с базовым, а итоговое ускорение равно 1,00.
При шаге 4 число прочитанных постингов на test сокращается
с 28\,832\,670 до 12\,635\,580, то есть примерно на 56,18\%,
но среднее время растёт с 33,326 до 82,122 мс. Таблица при этом
добавляет 4\,168\,634\,952 байта.

Контрольные суммы одинаковы для всех шагов отдельно на train и test.
Счётчики чтения в таблице суммируются по 15 измеряемым повторам.
Медиана на train при шаге 64 составляет 27,746 мс против 63,181 мс
без прыжков, однако среднее ухудшается с 71,839 до 175,074 мс.
Это показывает неоднородность запросов: улучшение типичного запроса
не гарантирует улучшения среднего по всему набору.

\conclusions
Механизм прыжков корректно сокращает чтение постингов, но в этой
реализации не дал выигрыша по выбранному критерию. Вероятные причины
- поиск по таблице при частых вызовах advance, случайные обращения
и большой служебный объём. Они требуют профилирования, а не следуют
однозначно из одного замера. Возможные улучшения: хранение таблиц
в памяти, кэширование подходящего блока между вызовами,
адаптивный шаг и отключение прыжков для коротких списков.
